\chapter{À LA RECHERCHE DE LA LANGUE ORIGINELLE}
\origpage{16}{002}
\begingroup\small\setstretch{1}\setlist[enumerate]{topsep=1pt,itemsep=1pt,parsep=0pt}\setlength{\parskip}{2pt}
\minititle{I.\quad LES RÉCITS ÉTIOLOGIQUES}
\begin{enumerate}[label=\arabic*.]
\item La tour de Babel : l’origine des langues
\item Langues sémitiques
\end{enumerate}
\minititle{II.\quad LES RÉCITS BIBLIQUES}
\begin{enumerate}[label=\arabic*.]
\item « Au commencement était le Verbe »
\item « Dieu dit : que la lumière soit, et la lumière fut. »
\item La Pentecôte : le don des langues
\begin{enumerate}[label=\textbf{\arabic*},nosep]
\item La xénoglossie
\item La glossolalie
\end{enumerate}
\end{enumerate}
\minititle{III.\quad LA RÉPONSE DE LA PHILOSOPHIE}
\begin{enumerate}[label=\arabic*.]
\item Aristote : « L’homme est un animal politique. »
\item Rousseau : « L’invention de la parole vient des passions. »
\end{enumerate}
\minititle{IV.\quad LA RÉPONSE DE LA SCIENCE}
\begin{enumerate}[label=\arabic*.]
\item Le point de vue de la paléontologie
\begin{enumerate}[label=\textbf{\arabic*},nosep]
\item \textit{L’Homo Habilis}
\item \textit{L’Homo Erectus}
\item \textit{L’Homo Sapiens}
\end{enumerate}
\item Le point de vue de la génétique
\begin{enumerate}[label=\textbf{\arabic*},nosep]
\item Diversité génétique et diversité linguistique
\item Les limites de l’analogie
\end{enumerate}
\item Le point de vue de la paléogénétique
\begin{enumerate}[label=\textbf{\arabic*},nosep]
\item La thèse de Philip Lieberman
\item L’homme de Néandertal : un grand bavard ?
\end{enumerate}
\end{enumerate}
\origpage{17}{003}
\minititle{Exercices}
\minititle{Pour aller plus loin}
\minititle{Glossaire}
\endgroup
\clearpage
\origpage{18}{004}
\lecturetitle{Introduction}
Fondée en 1864, la Société de Linguistique de Paris, a vu passer d’illustres linguistes français comme Michel Bréal, Antoine Meillet, Émile Benveniste ou encore André Martinet. Deux ans après sa création, une révision de ses statuts a exclu toute communication relative l’origine du langage,\ednote{原文 « relative l’origine » 缺介词，应作 « relative à l’origine »。} comme l’indiquait l’article 2 :

« La société n’admet aucune communication concernant soit l’origine du langage, soit la création d’une langue universelle. »

Lorsqu’il est arrivé à Paris en 1880, le linguiste suisse Ferdinand de Saussure, le père de la linguistique moderne, a fréquenté cette société. Dans son célèbre \textit{Cours de linguistique générale} (1916), il balaye d’un revers de main la question de l’origine du langage :

« La question de l’origine des langues n’a pas l’importance qu’on lui donne. Cette question n’existe même pas. Question de la source du Rhône : puérile ! Le moment de la genèse n’est lui-même pas saisissable : on ne le voit pas. »

Dans un séminaire de 1972, le psychanalyste Jacques Lacan (1901-1981) a lancé une prohibition semblable à celle de Saussure, pour que la linguistique soit digne d’être une science.

Cette interdiction a depuis peu été levée, nous voici donc libres de nous interroger sur la question de l’origine du langage, mais aussi des langues. Le langage est la capacité, propre à l’espèce humaine, de communiquer au moyen d’un système de signes vocaux (et, par-delà, au moyens de signes écrits).\ednote{原文 « au moyens » 应为 « au moyen »；固定结构为 « au moyen de »。} La langue est le langage utilisé par une communauté linguistique.

Pour tenter de répondre à cette double question, nous ferons appel aux mythes, à la religion, la philosophie, et la science.

\section{LES RÉCITS ÉTIOLOGIQUES}
L’étiologie vient d’un mot grec qui signifie « recherche des causes ». En médecine, l’étiologie désigne la recherche des causes d’une maladie. En philosophie et en littérature, un récit étiologique est une histoire qui a pour but de donner une explication imagée à un l’origine d’un phénomène.\ednote{原文多出 « un »；应为 « à l’origine d’un phénomène »。同段末句 « Un célèbre récit étiologique concerne le langage est… » 缺少关系代词 « qui »，应作 « Un célèbre récit étiologique qui concerne le langage est… »。} Ces récits expliquent l’origine du monde, des paysages, de l’homme, des animaux, des plantes, mais aussi du langage. Un célèbre récit étiologique concerne le langage est celui de la tour de Babel.

\subsection{La tour de Babel : l’origine des langues}
L’histoire de la tour de Babel est un épisode biblique rapporté dans la Genèse. Alors qu’ils parlent tous la même langue, les hommes atteignent une plaine dans le pays de Shinar et s’y installent tous. Là, ils entreprennent par eux-mêmes de bâtir une ville et une tour dont le sommet touche le ciel, pour se faire \origcont{19}{005}un nom. Alors Dieu brouille leur langue afin qu’ils ne se comprennent plus, et les disperse sur toute la surface de la Terre. La construction cesse. La ville est alors nommée Babel. Voici le texte original :

\begin{lecture}{Lecture : la Genèse (extrait)}
[1] Tout le monde se servait d’une même langue et des mêmes mots.\par
[2] Comme les hommes se déplaçaient à l’orient, ils trouvèrent une vallée au pays de Shinéar et ils s’y établirent.\par
[3] Ils se dirent l’un à l’autre : Allons ! Faisons des briques et cuisons-les au feu ! La brique leur servit de pierre et le bitume leur servit de mortier.\par
[4] Ils dirent : Allons ! Bâtissons-nous une ville et une tour dont le sommet pénètre les cieux ! Faisons-nous un nom et ne soyons pas dispersés sur toute la terre !\par
[5] Or Yahweh descendit pour voir la ville et la tour que les hommes avaient bâties.\par
[6] Et Yahweh dit : Voici que tous font un seul peuple et parlent une seule langue, et tel est le début de leurs entreprises ! Maintenant, aucun dessein ne sera irréalisable pour eux.\par
[7] Allons ! Descendons ! Et là, confondons leur langage pour qu’ils ne s’entendent plus les uns les autres.\par
[8] Yahvé les dispersa de là sur toute la face de la terre et ils cessèrent de bâtir la ville.\par
[9] Aussi la nomma-t-on Babel, car c’est là que Yahvé confondit le langage de tous les habitants de la terre et c’est de là qu’il les dispersa sur toute la face de la terre.
\end{lecture}
Si l’on fait une lecture « en linguiste » des quelques lignes ci-dessus, on peut en extraire quelques idées sur lesquelles je souhaite attirer votre attention :
\begin{itemize}
\item « Tout le monde se servait d’une même langue et des mêmes mots. » Pendant de nombreux siècles, on a pensé que cette langue mère était l’hébreu. Au 18\textsuperscript{e} siècle, c’est le sanskrit, langue sacrée de l’Inde antique, que certains linguistes ont considéré comme la mère de toutes les langues. Ce dernier a été ultérieurement détrôné au profit d’une langue hypothétique reconstruite, le proto-indo-européen\origfoot{1}{Tout cela sera développé dans le chapitre 5 de notre livre : Linguistique comparative et typologique.}.
\item Deuxième remarque : selon le mythe de Babel, toutes les langues seraient nées dans une région unique, à savoir Babylone, en Mésopotamie. Ce n’est ni tout à fait vrai, ni tout à fait faux. Ce qui est né en Mésopotamie, de façon établie et indubitable, ce ne sont pas les langues, mais l’écriture, avec les écritures les plus anciennes que sont le sumérien et l’akkadien.\ednote{这里混用了语言和文字系统。苏美尔语、阿卡德语是语言，楔形文字是用来记录这些语言的文字系统；此外，“文字起源于美索不达米亚，确定且无可置疑”的表述过强。牛津大学 ETCSL 对最早有文字记录的语言同时列出苏美尔语和古埃及语为候选。\ecite{2}} En revanche, l’idée selon laquelle le langage serait né dans un endroit unique avant de se diffuser sur toute la terre a été confirmée par la paléoanthropologie : c’est la théorie de l’origine africaine de l’homme moderne, connue du grand public sous le nom anglais d’\textit{Out of Africa} (la sortie d’Afrique). Appliquée au langage, cette théorie de l’origine unique de l’espèce s’appelle la monogenèse des langues\origfoot{2}{Nous reviendrons en détail sur le rapport entre linguistique et évolutionnisme dans les chapitres 5 (Linguistique comparative et typologique) et 16 (Biolinguistique et psycholinguistique).}\ednote{两点应区分：现代人类的起源证据本身不能逻辑上证明所有语言起源于同一种语言；这里把“人群起源”直接推作“语言单一起源”的证成，推论过强。另，原注 2 的 « 16 (Biolinguistique et psycholinguistique) » 与本书总目录不符，应为第 15 章；第 16 章为 « Langage et cognition : la linguistique cognitive »。}.
\end{itemize}
\origpage{20}{006}
\begin{itemize}
\item « Voici que tous font un seul peuple et parlent une seule langue. » Nous verrons que pour un pays, se doter d’une langue nationale, c’est assurer l’unité, la cohésion et l’indivisibilité de la nation. C’est la raison pour laquelle au lendemain de la révolution française, on a tenté d’éradiquer les langues régionales afin de ne pas fragmenter la nation et de faire émerger le français comme langue nationale. De nos jours, les questions ayant trait au choix d’une langue officielle ou nationale, au statut des langues régionales sont des enjeux majeurs relevant de la politique linguistique de chaque pays\origfoot{3}{Nous reviendrons sur cela dans le chapitre 13 : Sociolinguistique et géolinguistique.}.
\end{itemize}
\subsection{Langues sémitiques}
Les langues anciennes de Mésopotamie (sumérien, phénicien, akkadien, araméen, etc.) appartiennent à la famille des langues sémitiques.\ednote{原文把苏美尔语列为闪米特语，系分类错误。苏美尔语是孤立语言，没有已被可靠确立的亲属语言；阿卡德语属于闪米特语。应把 « sumérien » 从这一闪米特语列举中移出，并另作说明。参见牛津大学 ETCSL 的 « Sumerian language »。\ecite{2}} L’alphabet phénicien a été l’un des tous premiers systèmes d’écriture, utilisé puis transmis par les marchands phéniciens dans le monde méditerranéen où il a évolué et a été assimilé par de nombreuses cultures. L’alphabet araméen, forme modifiée du phénicien, est l’ancêtre de l’alphabet arabe moderne. L’alphabet hébreu moderne est lui-aussi issu de l’araméen. L’alphabet grec, et par extension ses descendants (l’alphabet latin, cyrillique et copte) sont des descendants directs de l’alphabet phénicien.\ednote{此处 « directs » 与前半句的层级关系不一致：按本句本身给出的链条，拉丁、西里尔和科普特字母经由希腊字母这一中间层，不能一概称为腓尼基字母的“直接”后裔；宜区分直接与间接传承。}
\begin{bbcomplement}{Babel : entre mythe et réalité}
Si l’histoire de la colère de Dieu est un mythe, la ville de Babylone a bien existé et se trouve aujourd’hui en Irak, dans la province de Bâbil, à une centaine de kilomètres au sud de Bagdad. Babylone était le siège de deux pouvoirs. Le premier était celui du roi, qui habitait un gigantesque palais situé au nord de la ville, le second celui du dieu de la ville, Mardouk, appelé aussi Bêl, roi des dieux du panthéon babylonien. La tour de Babel était un édifice religieux édifié en son honneur. Si l’Égypte a ses pyramides, Babylone avait ses \textit{ziggurat}, motakkadien (langue de Mésopotamie) signifiant « construit en hauteur ».\ednote{原文 « motakkadien » 中漏空格，应为 « mot akkadien »。下一句 « Édifiée sous les règnes des rois » 是本书使用的历史叙述，不据此假定《创世记》故事可与考古建筑逐项等同。} Édifiée sous les règnes des rois Nabopolassar (625-604 avant J.-C.) et Nabuchodonosor II (604-561 avant J.-C.), elle mesurait sept étages, faisait 91 mètres de côté, et était nommée \textit{Etemenanki}, ce qui en sumérien (une autre langue de Mésopotamie) signifie « la maison du ciel et de la terre ». La tour a bien sûr aujourd’hui disparu, mais son emplacement exact est connu. Pourquoi trouve-t-on alors un monument babylonien bien réel dans un récit mythique ? C’est l’histoire qui nous donne la réponse. Lorsque le roi babylonien Nabuchodonosor II a conquis la Judée en 586 avant J.-C., une partie de la population juive a été déportée à Babylone, où elle a vu la tour. C’est lors de cet exil que rédacteurs du Pentateuque (les cinq premiers livres de la Bible) l’ont intégrée dans leur récit biblique.\ednote{原文 « que rédacteurs » 缺少冠词，应为 « que les rédacteurs »。}
\end{bbcomplement}
\section{LES RÉCITS BIBLIQUES}
Le récit étiologique de Babel est un bon point de départ, mais il a un gros défaut : il explique l’origine des
\origcont{21}{007}« langues », mais pas celle de cette même langue originelle que parlaient les habitants de Babel. Il n’explique non plus quand et comment cette langue est apparue. La réponse à cette question se trouve dans la Bible.

\subsection{« Au commencement était le Verbe »}
Le texte suivant est extrait de l’\textit{Évangile selon Saint Jean}, le dernier des quatre évangiles du \textit{Nouveau Testament}. Saint Jean, c’est Jean de Zébédée, le chouchou de Jésus, nommé plusieurs fois « le disciple que Jésus aimait », « le disciple bien-aimé ».
\begin{lecture}{Lecture : la Vulgate (extrait)}
Au commencement était le Verbe, et le Verbe était avec Dieu, et le Verbe était Dieu.\par
Il était au commencement avec Dieu.Toutes choses ont été faites par lui, et rien de ce qui a été fait n’a été fait sans lui.
\end{lecture}
« Le verbe était Dieu » signifie l’identité entre le Verbe et de Dieu lui-même : le langage, attribut divin, et de nature divine devient consubstantiel à Dieu. Ainsi, se poser la question de l’origine du langage revient à se poser la question de l’origine de Dieu. Le langage, comme Dieu, a toujours existé, puisqu’il est « Dieu ».

Voilà qui répond à la question du « comment » et du « quand ». La réponse, un peu simple, tient en un mot : Dieu.

La deuxième phrase de ce très court extrait, « toutes choses ont été faites par lui », nous rappelle que c’est avec son Verbe que Dieu a créé le monde, comme nous allons le voir tout de suite.

\subsection{« Dieu dit : que la lumière soit, et la lumière fut. »}
Le mot « bible » vient du grec ancien \textit{biblia}, qui signifie « les livres ». Ce n’est en effet pas un livre, mais un ensemble de livres écrits par divers auteurs à diverses époques. Le Pentateuque (du grec \textit{penta}, cinq et \textit{teukhos}, livre en forme de rouleau) désigne les cinq premiers livres de la Bible. Le premier des cinq livres du Pentateuque, la Genèse, raconte la création de l’univers et de l’homme.
\begin{lecture}{Lecture : la Genèse (extrait)}
[1] Au commencement, Dieu créa les cieux et la terre.\par
[2] La terre était informe et vide : il y avait des ténèbres à la surface de l’abîme, et l’esprit de Dieu se mouvait au-dessus des eaux.\par
[3] Dieu dit : Que la lumière soit ! Et la lumière fut.\par
[4] Dieu vit que la lumière était bonne ; et Dieu sépara la lumière d’avec les ténèbres.\par
[5] Dieu appela la lumière jour, et il appela les ténèbres nuit. Ainsi, il y eut un soir, et il y eut un matin : ce fut le premier jour [...]
\end{lecture}
\origpage{22}{008}
Si l’on fait une lecture non pas biblique mais linguistique du texte ci-dessus, on peut y trouver un certain nombre de concepts importants que nous développerons plus tard dans notre livre. En voici deux exemples:
\begin{itemize}
\item « Dieu dit » exprime clairement que c’est le Verbe de Dieu, puissance créatrice, qui a créé le monde. Dans cette conception, le langage est vu comme une force agissante qui permet de modifier son environnement. La fonction du langage n’est pas seulement d’exprimer nos pensées ou nos émotions. Il permet aussi de « faire des choses ». Lorsque dans le mythe de Babel Dieu confond les langues, c’est pour empêcher les hommes « faire » quelque chose,\ednote{原文 « empêcher les hommes « faire » » 中缺少 « de »，应作 « empêcher les hommes de « faire » quelque chose »。} à savoir bâtir une tour. Le langage a pour finalité l’action : c’est le point de vue du philosophe français Henri Bergson, que nous présenterons un peu plus tard dans ce chapitre.
\item « Dieu appela la lumière jour, et il appela les ténèbres nuit. » Après avoir créé toutes les choses ici-bas, il fallait bien leur donner un nom. Ce processus est appelé « nomination ». Dieu appela la lumière « jour », et il appela les ténèbres « nuit ». Pourquoi pas. Mais aurait-il pu appeler la lumière autrement ? Aurait-il pu appeler la lumière « nuit » et les ténèbres « jour » ? Les noms donnés aux objets sont-ils exacts ? Désignent-ils l’essence de la chose nommée ? Le chapitre 4 (Logique et philosophie du langage) nous apportera des éléments de réponse.
\end{itemize}
\subsection{La Pentecôte : le don des langues}
Restons dans les récits bibliques traitant indirectement du langage. Le récit des \textit{Actes des Apôtres} est le cinquième livre du \textit{Nouveau Testament}. Ce récit débute avec l’\textit{Ascension}, qui marque l’élévation au ciel de Jésus Christ, et se poursuit avec la \textit{Pentecôte}, c’est-à-dire la venue du Saint Esprit sur cent-vingt disciples de Jésus. Le mot « Pentecôte »vient du grec \textit{pentaekoste hemerra}, le cinquantième jour, car la fête de la Pentecôte a lieu cinquante jours après Pâques.
\begin{lecture}{Lecture : \textit{Actes de apôtres} (extrait)\ednote{标题原文 « Actes de apôtres » 应为 « Actes des apôtres »；本页上文也使用 « Actes des Apôtres »。}}
[1] Le jour de la Pentecôte, ils étaient tous ensemble dans le même lieu.\par
[2] Tout à coup il vint du ciel un bruit comme celui d’un vent impétueux, et il remplit toute la maison où ils étaient assis.\par
[3] Des langues, semblables à des langues de feu, leur apparurent, séparées les unes des autres, et se posèrent sur chacun d’eux.\par
[4] Et ils furent tous remplis du Saint-Esprit, et se mirent à parler en d’autres langues, selon que l’Esprit leur donnait de s’exprimer.
\end{lecture}
La langue maternelle des apôtres était le galiléen, langue parlée dans la région de la Galilée, en Israël. Or, voilà que soudain, ils se mettent à parler d’autres langues, sans les avoir apprises ! On comprend d’où vient l’expression « avoir le don des langues ». Connaissant de nombreuses langues, les apôtres allaient pouvoir remplir leur mission itinérante et répandre la Parole de Dieu dans différentes régions du monde. La question que l’on doit se poser, c’est comment peut-on subitement parler des langues sans les avoir apprises ?
\origpage{23}{009}
Et pourtant, ce phénomène, bien que fort rare, a déjà été observé à plusieurs reprises dans l’histoire. Non seulement certaines personnes parlent des langues sans les avoir apprises, mais il s’agit en plus de langues mortes et connues seulement de rares spécialistes. Ce phénomène est désigné sous les noms de « xénoglossie » (du grec \textit{xéno}, étranger, et \textit{glossis}, parler) et de « glossolalie » (\textit{glôssa}, langue, et \textit{laléô}, parler).\ednote{此处两种术语的希腊语词源解释不一致：« glôssa/glōssa » 指“舌、语言”，不是“说话”；“说话”在下一词的解释中对应 « laléô »。因此 xénoglossie 的解释至少应把 « glossis, parler » 更正为指“语言”的成分。这里只纠正本段的词义混淆，超常语言能力的案例仍属待核，不据本段叙事当作已证实事实。}
\subsubsection*{1\quad La xénoglossie}
L’épisode du don des langues de la Pentecôte relate un cas de xénoglossie, c’est-à-dire le fait de savoir parler ou d’écrire (xénographie) une langue étrangère existante sans l’avoir apprise. Dans l’histoire de France, des phénomènes de xénoglossie se seraient manifestés parmi les protestants après la révocation par Louis XIV de l’Édit de Nantes : on raconte que des enfants illettrés exhortaient dans un français excellent leur entourage qui ne s’exprimait qu’en patois. Des paysans incultes se seraient mis à parler en latin, en grec et en hébreu. Sur ces récits anciens, on peut avoir quelques doutes, mais les observations récentes ont de quoi stupéfier (et faire envie). Par exemple, dans les années 1930, les deux jumeaux communiquaient entre deux dans une langue étrangère que leur propre père, le médecin new-yorkais Marshall Duffie, ne comprenait pas. Les professeurs de langues étrangères de l’Université de Columbia ont été eux aussi incapables de dire de quelle langue il s’agissait. C’est finalement un professeur de langues anciennes qui a découvert que les deux enfants parlaient l’araméen, c’est-à-dire « la langue de Jésus-Christ ». En Allemagne, le cas avéré de la mystique catholique Thérèse Neumann (1898-1962) qui parlait en araméen sans l’avoir appris, est aussi très célèbre.
\subsubsection*{2\quad La glossolalie}
Dans le cas de la glossolalie, la langue parlée est fictive, mais elle est intelligible en ce sens qu’elle possède une structure grammaticale et un vocabulaire constant. Au début du siècle dernier, la médium suisse Hélène Smith qui prétendait pouvoir parler avec les morts (notamment Victor Hugo), parlait et écrivait en sanskrit en araméen, mais aussi en langage martien. Le psychologue suisse Théodore Flournoy (1854-1920,\ednote{原文生卒年后的右括号缺失，应为 « Théodore Flournoy (1854-1920), spécialiste… »。} spécialiste du spiritisme et du paranormal, mais aussi le grand linguiste Ferdinand de Saussure se sont intéressés à son cas. Comment expliquer ces phénomènes ? Le psychiatre Ian Stevenson évoque deux raisons. La première, c’est la réincarnation\origfoot{4}{Réincarnation : dans certaines religions (hindouisme, bouddhisme, etc.), croyance selon laquelle l’âme au moment de la mort passe dans un autre corps.}, de langues parlées dans une autre vie. Hélène Smith, lorsqu’elle parlait le sanskrit, était la réincarnation d’une princesse indienne. L’autre raison avancée, un peu plus crédible pour un esprit rationnel, est celle de la réminiscence, c’est-à-dire des langues auxquelles ces personnes ont été confrontées dans l’enfance, avant de les oublier à l’âge adulte. C’est ce que Stevenson appelle la xénoglossie récitative. Pour d’autres, ces phénomènes relèvent du miracle divin, ou de la maladie mentale.

Les récits mythologiques et bibliques que nous avons lus jusqu’à présent ont l’avantage d’être facilement compréhensibles, parce qu’imagés, mais on ne peut les prendre pour argent comptant. Réfléchir à la question de l’origine des langues et du langage en se basant sur une pensée non \origcont{24}{010}critique, qui accepte sans remise en question les croyances et les mythes, cela n’est pas possible. Il faut rompre avec la vision mythologique du monde et passer à l’étape suivante, la philosophie.
\section{LA RÉPONSE DE LA PHILOSOPHIE}
Pour Aristote, « la philosophie commence avec l’étonnement ». Partant de l’ignorance et de l’étonnement des hommes face aux choses du monde, il justifie la naissance de la philosophie qui n’a pas d’autre but que de satisfaire la soif de savoir des hommes au moyen de la raison humaine (\textit{logos}, en grec). Dans \textit{La Politique}, et dans l’extrait ci-dessous, Aristote relie l’invention du langage aux valeurs de bien et de justice. Plus près de nous, durant le siècle des lumières, Rousseau s’est lui aussi intéressé au langage : nous présenterons dans cette troisième partie de larges extraits de son \textit{Essai sur l’origine des langues}. Mais à tout seigneur, tout honneur : commençons par Aristote.
\subsection{Aristote : « L’homme est un animal politique. »}
\begin{lecture}{Lecture : \textit{La Politique} (extrait)}
La cité est au nombre des réalités qui existent naturellement, et [...] l’homme est par nature un animal politique. [...] Mais que l’homme soit un animal politique à un plus haut degré qu’une abeille quelconque ou tout autre animal vivant à l’état grégaire, cela est évident. La nature, en effet, selon nous, ne fait rien en vain ; et l’homme seul de tous les animaux, possède la parole. Or, tandis que la voix ne sert qu’à indiquer la joie et la peine, et appartient aux animaux également (car leur nature va jusqu’à éprouver les sensations de plaisir et de douleur, et à se les signifier les uns aux autres), le discours sert à exprimer l’utile et le nuisible, et par suit aussi,\ednote{底本引文作 « par suit aussi »，此处固定短语应为 « par suite aussi »。} le juste et l’injuste : car c’est le caractère propre à l’homme par rapport aux autres animaux, d’être le seul à avoir le sentiment du bien et du mal, du juste et de l’injuste, et des autres notions morales, et c’est la communauté de ces sentiments qui engendre famille et cité.
\end{lecture}
Le texte ci-dessus repose sur deux caractéristiques humaines étroitement liées : parler et vivre en société. Si l’homme possède le langage, c’est parce qu’il est fait pour vivre en société. Voilà pour l’idée centrale. Ce qui est plus intéressant, c’est qu’Aristote nous demande de faire une distinction entre « voix » (en grec, \textit{phônê}) et la « parole » (en grec \textit{logos}, traduite aussi par « discours » dans le texte). La voix est expressive, et permet de communiquer ses affects (plaisir, peine, douleur, etc.). On la rencontre donc chez les animaux comme chez les hommes, l’homme étant un animal. Politique, certes, mais animal tout de même.

Or, l’homme, en plus de la « voix », dispose de « la parole », et cette dernière dépasse la fonction purement expressive et communicative de la « voix ». La \textit{phônê} concerne les êtres sentants, mais pour Aristote, seul un être pensant, doté de raison peut véritablement parler. C’est que les Grecs soulignent avec la notion de \textit{logos} qui signifie aussi « parole », « discours » mais aussi « raison ». Dans le contexte de la démocratie athénienne, c’est le « discours » qui permet de débattre des valeurs de la communauté, et cela est précisément l’enjeu de l’activité politique. Il s’ensuit que pour Aristote, le langage est d’essence politique et inversement, la politique est d’essence langagière. L’espace politique est le lieu où chacun peut exprimer sa conception de l’utile, du \origcont{25}{011}bien, du juste, de l’utile et du nuisible, autant de valeurs qui ne seraient pas exprimables au seul moyen de la « voix ». Le langage est l’instrument du débat, moyen par lequel une pluralité d’êtres différents peut aboutir à des accords nécessaires à la vie en commun dans la Cité.
\subsection{Rousseau : « L’invention de la parole vient des passions. »}
Si tout le monde connaît \textit{Émile} ou \textit{De l’éducation} (1762), le \textit{Contrat social} (1762), de \textbf{Jean-Jacques Rousseau (1712-1778)}, son \textit{Essai sur l’origine des langues} (\textit{Essai sur l’origine des langues où il est parlé de la mélodie et de l’imitation musicale}) n’a pas connu la même renommée. Dans cet essai publié en 1781, Rousseau complète son célèbre \textit{Discours sur l’origine et les fondements de l’inégalité parmi les hommes} de 1755.

Contrairement aux explications simples et quelques peu naïves des mythes et récits bibliques, Rousseau aborde la question de l’origine des langues avec la rationalité des philosophes des Lumières.
\begin{lecture}{Lecture : \textit{Essai sur l’origine des langues} (extrait 1)}
La parole distingue l’homme entre les animaux : le langage distingue les nations entre elles ; on ne connaît d’où est un homme qu’après qu’il a parlé. L’usage et le besoin font apprendre à chacun la langue de son pays[...] La parole étant la première institution sociale ne doit sa forme qu’à des causes naturelles. Sitôt qu’un homme fut reconnu par un autre pour un être sentant, pensant et semblable à lui, le désir ou le besoin de lui communiquer ses sentiments et ses pensées lui en fit chercher les moyens. Ces moyens ne peuvent se tirer que des sens, les seuls instruments par lesquels un homme puisse agir sur un autre. Voilà donc l’institution des signes sensibles pour exprimer la pensée.
\end{lecture}
Pour guider un peu votre lecture, j’attire votre attention sur les trois notions suivantes :
\begin{center}\small
\renewcommand{\arraystretch}{1.2}
\begin{tabularx}{\linewidth}{|>{\RaggedRight\arraybackslash}X|>{\RaggedRight\arraybackslash}X|>{\RaggedRight\arraybackslash}X|}
\hline
\centering\textbf{LANGAGE}&\centering\textbf{LANGUE}&\centering\arraybackslash\textbf{PAROLE}\\\hline
Aptitude à communiquer, propre à l’être humain, au moyen de signes vocaux et graphiques correspondant à un sens.&Produit acquis : instrument de communication ; code constitué en un système de règles communes à une même communauté.&Utilisation individuelle de la langue par un sujet parlant.\\\hline
\centering\textbf{Langage/langue}&\centering\textbf{La langue/une langue}&\centering\arraybackslash\textbf{Langue/parole}\\\hline
\centering\textbf{Abstrait/concret}&\centering\textbf{Général/particulier}&\centering\arraybackslash\textbf{Tout/partie}\newline\textbf{Social/individuel}\\\hline
\end{tabularx}
\end{center}
Retenez bien que le langage est une faculté abstraite et universelle, alors que la langue, qui est concrète et non innée, concerne une communauté ou un groupe linguistique. Quant à la parole, elle se rapporte à un locuteur.

En gardant ces éléments en tête, nous allons commenter l’extrait de Rousseau présenté un peu plus haut :
\origpage{26}{012}
\begin{itemize}
\item « Le langage distingue les nations entre elles. » Le langage est une faculté universelle, identique chez tous les hommes de tous peuples et de toutes nations. Ce qui distingue les nations entre elles, ce n’est pas \textit{le langage}, ce sont \textit{les langues}.
\item « L’usage et le besoin font apprendre à chacun la langue de son pays. » D’après Rousseau,les langues, elles, sont acquises et doivent faire l’objet d’un apprentissage. Le linguiste américain \textbf{Noam Chomsky} et psychologue \textbf{Steven Pinker} soutiennent l’opinion contraire selon laquelle les langues ne sont pas acquises mais innées\origfoot{5}{Cette question qui touche à l’acquisition des langues a suscité de vifs débats entre linguistes et psychologues. Nous étudierons cela en détail dans le chapitre 15, consacré à la biolinguistique et à la psycholinguistique.}\ednote{此句把“语言能力的先天基础”与“某种具体语言无需习得”混为一谈。应改为：他们强调语言能力或语言习得机制具有先天基础，而不是认为法语、汉语等具体语言不经接触即可获得。Pinker 对《语言本能》的作者说明明确同时讨论语言的本能基础和儿童如何习得语言。\ecite{3}}.
\item « La parole étant la première institution sociale. » Ici, Rousseau a tort et raison à la fois. Sur la dimension sociale du langage, il reprend l’idée d’Aristote : l’homme est un animal politique fait pour vivre en société, et le langage est la condition de cette vie sociale. Soit. Mais la « parole », au sens de langage n’est pas une institution : elle est une faculté universelle et naturelle. L’institution, ce sont les langues particulières et sociales. Les hommes du monde entier partagent une même faculté de langage, mais parlent des langues différentes. Les locuteurs de ces langues appartiennent à une même communauté linguistique, c’est en cela que les langues ont une dimension sociale. Dans son \textit{Cours de linguistique générale} (1916), \textbf{Ferdinand de Saussure}, en différenciant langue et « parole » va donner à cette dernière le sens que nous lui connaissons aujourd’hui, à savoir l’utilisation concrète par un individu et dans une situation particulière de cette faculté abstraite et universelle qu’est le langage\origfoot{6}{Nous reviendrons longuement sur tout cela dans le chapitre 6 (Saussure, et la naissance de la linguistique moderne).}.
\item « Un être pensant et sentant » : Rousseau fait ici référence à Blaise Pascal qui a écrit : « L’homme est un roseau, le plus faible de la nature, mais c’est un roseau pensant. » Nous ferons aussi le lien avec le texte d’Aristote : l’être « sentant » n’a besoin que de la « voix » (\textit{phônê}), l’être « pensant » a besoin du langage (\textit{logos}).
\end{itemize}
En gardant bien sûr en tête ce que nous venons d’apprendre, continuons à présent avec le deuxième extrait.
\begin{lecture}{Lecture : \textit{Essai sur l’origine des langues} (extrait 2)}
Il paraît encore que l’invention de l’art de communiquer nos idées dépend moins des organes qui nous servent à cette communication, que d’une faculté propre à l’homme, qui lui fait employer ses organes à cet usage, et qui, si ceux-là lui manquaient, lui en ferait employer d’autres à la même fin.
\end{lecture}
Pourvu seulement qu’il y ait entre lui et ses semblables quelque moyen de communication par lequel l’un puisse agir et l’autre sentir ils parviendront à se communiquer enfin tout autant d’idées qu’ils en auront. Les animaux ont pour cette communication une organisation plus que suffisante, et jamais aucun d’eux n’en a fait cet usage. Voilà ce me semble, une différence bien caractéristique. Ceux d’entre eux qui travaillent et vivent en commun, les castors, \origcont{27}{013}les fourmis, les abeilles, ont quelque langue naturelle pour s’entre-communiquer je n’en fais aucun doute. Il y a même lieu de croire que la langue des castors et celle des fourmis sont dans le geste et parlent seulement aux yeux. Quoi qu’il en soit, par cela même que les unes et les autres de ces langues sont naturelles, elles ne sont pas acquises : les animaux qui les parlent les ont en naissant : ils les ont tous, et partout la même : ils n’en changent point, ils n’y font pas le moindre progrès. La langue de convention n’appartient qu’à l’homme.

Voici quelques idées sur lesquelles je souhaite attirer votre attention :
\begin{itemize}
\item « \textit{L’invention} de l’art de communiquer nos pensées » est à rapprocher de « \textit{l’institution} des signes » qui concluait le premier extrait. « invention » et « institution » expriment toutes deux l’idée que les hommes sont les artisans de leurs langues.
\item Lorsqu’il évoque « les organes » qui servent à la communication, Rousseau évoque ce que nous appelons de nos jours les organes de la phonation\origfoot{7}{Voir le chapitre 7 : Phonétique articulatoire.}. L’idée que Rousseau développe ici est tout à fait exacte : pour parler, nous utilisons des organes (poumons, bouche, nez, dents, langue, etc.) qui ne sont initialement pas destinés au langage. Leur fonction principale est la respiration et la mastication. L’usage de ces organes dans la phonation n’est que secondaire. L’homme n’a pas à proprement parler d’organe spécialement dédié du langage.
\item En cas d’absence ou de déficience des organes de la phonation (« si ceux-là lui manquaient »), l’homme peut communiquer sans eux : la langue des signes des sourds et muets en la preuve. La faculté langagière de l’homme ne se réduit pas au bon fonctionnement organes\origfoot{8}{Quant aux pathologies du langage, elles seront présentées dans notre chapitre 15 : Psycholinguistique.}\ednote{本条两处漏词：« en la preuve » 应为 « en est la preuve »；« au bon fonctionnement organes » 应为 « au bon fonctionnement des organes »。上一条 « dédié du langage » 应为 « dédié au langage »。}.
\item « Les animaux ont pour cette communication une organisation plus que suffisante, et jamais aucun d’eux n’en a fait cet usage. » Intuition très intéressante de Rousseau : les animaux disposent des mêmes moyens que nous, c’est-à-dire du même appareil phonatoire (poumons, bouche, nez, langue, etc.) mais ils ne s’en servent pas. Autrement dit, ils \textit{pourraient} parler, mais ils ne le font pas. Pour quelle raison ? Réponse dans le prochain chapitre, avec le philosophe René Descartes.
\item « La langue de convention n’appartient qu’à l’homme. » Nous avons soulevé un peu plus haut dans ce chapitre la question du rapport qui entre un mot et l’objet qu’il désigne. Rousseau défend ici le point de vue conventionnel, c’est-à-dire que les mots de notre langue ont d’abord été inventés et assignés aux objets les hommes eux-mêmes, avant d’être acceptés (d’où le mot de convention) comme tels par la communauté linguistique. Si l’objet que vous tenez dans votre main s’appelle « livre », c’est parce qu’il y a fort longtemps les hommes se sont mis d’accord pour nommer ainsi cet objet. C’est la première convention. Et aujourd’hui, vous utilisez le mot « livre » lorsque vous avez commencé à apprendre le français, votre professeur vous a dit : « 法语中，livre 就是书的意思 », et vous avez accepté cela. C’est la deuxième convention.
\end{itemize}
Pour terminer sur ce deuxième extrait, je vous incite à relever l’usage répété des marque de réserve et de prudence de la part de Rousseau (« il paraît », «ce me semble», « il y a lieu de croire »). C’est d’ailleurs sur une \origcont{28}{014}de ces marques de prudence que s’ouvre notre troisième extrait.
\begin{lecture}{Lecture : \textit{Essai sur l’origine des langues} (extrait 3)}
Il est donc à croire que les besoins dictèrent les premiers gestes, et que les passions arrachèrent les premières voix. Et, suivant avec ces distinctions la trace des faits, peut-être faudrait-il raisonner sur l’origine des langues tout autrement qu’on a fait jusqu’ici. [...] On prétend que les hommes inventèrent la parole pour exprimer leurs besoins ; cette opinion me paraît insoutenable. L’effet naturel du besoin fut d’écarter les hommes et non de les rapprocher. Il le fallait ainsi pour que l’espèce vînt à s’étendre et que la terre se peuplât promptement, sans quoi le genre humain se fût entassé dans un coin du monde, et tout le reste fût demeuré désert. De cela seul il suit que l’origine des langues n’est point due aux premiers besoins des hommes ; il serait absurde que de la cause qui les écarte vînt le moyen qui les unit. D’où peut donc venir cette origine ? Des besoins moraux, des passions. Toutes les passions rapprochent les hommes que la nécessité de chercher à vivre force à se fuir. Ce n’est ni la faim ni la soif mais l’amour, la haine, la pitié, la colère qui leur ont arraché les premières voix. Les fruits ne se dérobent point à nos mains, on peut s’en nourrir sans parler, on poursuit en silence la proie dont on veut se repaître ; mais pour émouvoir un jeune cœur, pour repousser un agresseur injuste, la nature dicte des accents, des cris, des plaintes : voilà les anciens mots inventés, et voilà pourquoi les premières langues furent chantantes et passionnées avant d’être simples et méthodiques.
\end{lecture}
Cet extrait est lui aussi particulièrement marqué par la prudence : « il est donc à croire que », « peut-être », « me paraît ». La raison est que Rousseau s’apprête à formuler son hypothèse sur l’origine du langage. La prudence s’impose donc, d’autant que contrairement à ce que nous avons vu dans les récits étiologiques et bibliques, Dieu est absent du schéma explicatif de Rousseau. Comme précédemment, j’attire votre attention sur les passages importants :
\begin{itemize}
\item « Les besoins dictèrent les premiers gestes », « les passions arrachèrent les premières voix ». Ici, Rousseau semble opposer d’emblée « besoins » et « passions », mais ce n’est pas tout à fait le cas. Il va préciser plus bas dans le texte que ce qu’il appelle « passions » sont les « besoins moraux ». Il oppose donc en fait deux types de besoins différents : les besoins du corps, et les besoins moraux. Ce sont les besoins moraux qui ont rendu le langage nécessaire.
\item « On ne commença pas par raisonner mais par sentir. » Pour utiliser des adjectifs rencontrés plus tôt, l’homme, selon Rousseau, a d’abord été un être « sentant » avant de devenir un être « pensant ». De façon implicite, on retrouve l’idée d’Aristote d’une progression d’un niveau initial (la voix) à un niveau supérieur (le langage), ce dernier étant utilisé pour « raisonner » et exprimer des « besoins moraux », comme par exemple « émouvoir un jeune cœur » ou « repousser un agresseur injuste ». Souvenez-vous d’ailleurs que dans ce même texte d’Aristote, la nécessité du langage était déjà liée au besoin d’exprimer « le juste et d’injuste », c’est-à-dire des valeurs morales.\ednote{原文 « le juste et d’injuste » 应为 « le juste et l’injuste »，可与前引 Aristote 文本对校。}
\item « L’effet naturel du besoin fut d’écarter les hommes et non de les rapprocher. » Pour Rousseau, l’homme n’est pas un animal politique. Dans son \textit{Discours sur l’origine des inégalités parmi les \origcont{29}{015}hommes}, il exprime l’idée que « l’homme est naturellement bon [...], la société déprave et pervertit les hommes », l’idée qui vous est sûrement familière ici en Chine : « 人之初性本善，性相近习相远 ». « Naturellement », donc, les hommes souhaitent s’éloigner les uns des autres, « se fuir ». Rousseau réfute donc la conception d’Aristote vue plus haut selon laquelle l’origine du langage tient au fait que l’homme est fait par la cité.
\item « Des besoins moraux, des passions. » Ce sont les besoins « moraux » qui ont rendu le langage nécessaire. L’adjectif « moral » est important car il refait le lien avec une idée similaire chez :\ednote{原文 « chez : » 后缺少所指人名。下文引用与本节前引 Aristote 文本对应，故宜补为 « chez Aristote : »。} « C’est le caractère propre à l’homme par rapport aux autres animaux, d’être le seul à avoir le sentiment du bien et du mal, du juste et de l’injuste, et des autres notions morales. » Pour survivre, pour satisfaire ses besoins physiques, l’homme n’avait pas besoin de la parole (« Ce n’est ni la faim ni la soif [...] qui leur ont arraché les premières voix. »). C’est pour exprimer amour, haine, pitié, colère, indignation face à l’injustice que l’homme a inventé le langage.
\end{itemize}
\begin{bbcomplement}{Schopenhauer, \textit{Parerga et Paralipomena} (extrait)}
Par une froide journée d’hiver un troupeau de porcs-épics s’était mis en groupe serré pour se garantir mutuellement contre la gelée par leur propre chaleur. Mais tout aussitôt ils ressentirent les atteintes de leurs piquants, ce qui les fit s’écarter les uns des autres. Quand le besoin de se réchauffer les eut rapprochés de nouveau, le même inconvénient se renouvela, de sorte qu’ils étaient ballottés de çà et de là entre les deux maux jusqu’à ce qu’ils eussent fini par trouver une distance moyenne qui leur rendît la situation supportable. Ainsi, le besoin de société, né du vide et de la monotonie de leur vie intérieure, pousse les hommes les uns vers les autres ; mais leurs nombreuses manières d’être antipathiques et leurs insupportables défauts les dispersent de nouveau. La distance moyenne qu’ils finissent par découvrir et à laquelle la vie en commun.\ednote{底本引文末句在 « et à laquelle la vie en commun » 处缺文，句法不完整。本版保留断残原文；缺失文字尚未用对应法译版本确证，不凭意译补齐。}

\begin{center}\bfseries Rousseau, \textit{Discours sur l’origine de l’inégalité parmi les hommes}\end{center}
Qu’on songe de combien d’idées nous sommes redevables à l’usage de la parole ; combien la grammaire exerce et facilite les opérations de l’esprit ; et qu’on pense aux peines inconcevables, et au temps infini qu’a dû coûter la première invention des langues ; qu’on joigne ces réflexions aux précédents, et l’on jugera combien il eût fallu de milliers de siècles, pour développer successivement dans l’esprit humain les opérations dont il était capable.

Qu’il me soit permis de considérer un instant les embarras de l’origine des langues. Je pourrais me contenter de citer ou de répéter ici les recherches que M. l’Abbé de Condillac a faites sur cette matière, qui toutes confirment pleinement mon sentiment, et qui, peut-être, m’en ont donné la première idée. Mais la manière dont ce philosophe résout les difficultés qu’il se fait à lui-même sur l’origine des signes institués, \origcont{30}{016}montrant qu’il a supposé ce que je mets en question, savoir une sorte de société déjà établie entre les inventeurs du langage, je crois en renvoyant à ses réflexions devoir y joindre les miennes pour exposer les mêmes difficultés dans le jour qui convient à mon sujet.

Ayant eu des enfants et ceux-ci développant dans leurs échanges avec leurs parents une dextérité particulière de la langue, le langage articulé se développe au détriment du langage d’action parce qu’il se révèle plus « économique ». Selon lui, le langage est donc d’invention purement humaine, les signes et des langues sont « d’institution », et non de nature. La première qui se présente est d’imaginer comment elles purent devenir nécessaires ; car les hommes n’ayant nulle correspondance entre eux, ni aucun besoin d’en avoir, on ne conçoit ni la nécessité de cette invention, ni sa possibilité, si elle ne fut pas indispensable. Je dirais bien, comme beaucoup d’autres, que les langues sont nées dans le commerce domestique des pères, des mères et des enfants [...] cela montre bien comment on enseigne des langues déjà formées, mais cela n’apprend point comment elles se forment. Supposons cette première difficulté vaincue : franchissons pour un moment l’espace immense qui dut se trouver entre le pur état de nature et le besoin des langues ; et cherchons, en les supposant nécessaires, comment elles purent commencer à s’établir. Nouvelle difficulté pire encore que la précédente ; car si les hommes ont eu besoin de la parole pour apprendre à penser, ils ont eu bien plus besoin encore de savoir penser pour trouver l’art de la parole ; et quand on comprendrait comment les sons de la voix ont été pris pour les interprètes conventionnels de nos idées, il resterait toujours à savoir quels ont pu être les interprètes mêmes de cette convention pour les idées qui, n’ayant point un objet sensible, ne pouvaient s’indiquer ni par le geste, ni par la voix, de sorte qu’à peine peut-on former des conjectures supportables sur la naissance de cet art de communiquer ses pensées, et d’établir un commerce entre les esprits.

Le premier langage de l’homme, le langage le plus universel, le plus énergique, et le seul dont il eut besoin, avant qu’il fallût persuader des hommes assemblés, est le cri de la nature. Comme ce cri n’était arraché que par une sorte d’instinct dans les occasions pressantes, pour implorer du secours dans les grands dangers, ou du soulagement dans les maux violents, il n’était pas d’un grand usage dans le cours ordinaire de la vie, où règnent des sentiments plus modérés.

Quand les idées des hommes commencèrent à s’étendre et à se multiplier, et qu’il s’établit entre eux une communication plus étroite, ils cherchèrent des signes plus nombreux et un langage plus étendu : ils multiplièrent les inflexions de la voix, et y joignirent les gestes, qui, par leur nature, sont plus expressifs, et dont le sens dépend moins d’une détermination antérieure.

Ils exprimaient donc les objets visibles par des gestes, et ceux qui frappent l’ouïe, par des sons imitatifs mais comme le geste n’indique guère que les objets présents, ou faciles à décrire, et les actions visibles ; qu’il n’est pas d’un usage universel, puisque l’obscurité, ou l’interposition d’un corps le rend inutile, et qu’il exige l’attention plutôt qu’il ne l’excite, on s’avisa enfin de lui substituer les articulations de la voix, qui, sans avoir le même rapport avec certaines idées, sont plus propres à les représenter toutes, comme
\origcont{31}{017}signes institués ; substitution qui ne put se faire que d’un commun consentement, et d’une manière assez difficile à pratiquer pour des hommes dont les organes grossiers n’avaient encore aucun exercice, et plus difficile encore à concevoir en elle-même, puisque cet accord unanime dut être motivé, et que la parole paraît avoir été fort nécessaire, pour établir l’usage de la parole.
\end{bbcomplement}
Quittons à présent la philosophie pour entrer de plein pied dans le domaine de la science.\ednote{原文 « de plein pied » 应为 « de plain-pied »，这里意为直接进入。}
\section{LA RÉPONSE DE LA SCIENCE}
Grâce à la paléontologie et à la génétique, nous serons en mesure d’apporter des éléments de réponse factuels et crédibles à la question de l’origine du langage.
\subsection{Le point de vue de la paléontologie}
La paléontologie est la discipline scientifique qui étudie les restes fossiles des êtres vivants du passé et les implications évolutives ressortant de l’étude de ces restes. L’évolution de l’homme est assez souvent comparée visuellement à un buisson ou à un arbre. Une fois notre lignée séparée des grands singes il y a 7 millions d’années, plusieurs branches avec de nouvelles espèces se sont écartées, se sont rapprochées, ont disparu brusquement, pour finir avec une seule espèce depuis quelques milliers d’années : l’\textit{Homo Sapiens}.
\subsubsection*{1\quad L’\textit{Homo Habilis}}
L’\textit{Homo Habilis}, l’homme « adroit de ses mains », a vécu en Afrique entre environ 2, 3 et 1, 5 millions d’années avant notre ère. Découvert en 1961 au Nord de la Tanzanie, il pourrait avoir été le plus ancien préhumain à avoir employé une forme primitive de langage. Cette hypothèse se base sur les éléments objectifs suivants :
\begin{itemize}
\item Le redressement du crâne chez l’Homo Habilis a permis d’abaisser le pharynx et surtout le larynx, l’organe de la phonation. Cela permet la production d’un langage articulé et modulé.
\item L’abaissement des voies aériennes supérieures a permis en outre d’augmenter la hauteur de la voûte du palais, permettant ainsi à la langue d’articuler une plus large gamme de sons.
\item Enfin, les découvertes récentes des neurosciences ont montré que les zones du cerveau responsables du langage sont proches de celles mobilisées pour le travail manuel de précision. Or, l’Homo Habilis, contemporain de la pierre taillée, est l’inventeur de l’outil. Cela induit un développement simultané des facultés langagières et des facultés manufacturières du genre humain.
\end{itemize}
Sur les éléments ci-dessus, il n’y a pas de consensus, et les choses ne sont pas si simples. En 1983, \textbf{Ralf Holloway} a provoqué une secousse en annonçant qu’il avait repéré sur un crâne d’Homo Habilis la présence embryonnaire d’une zone cérébrale spécialisée dans le langage.
\origpage{32}{018}
Mais d’autres études tendent à démontrer que son larynx n’était pas descendu et qu’il ne pouvait donc pas articuler de manière correcte. C’est la thèse défendue par \textbf{Philip Tobias} (1925-2012) paléoanthropologue sud-africain et codécouvreur de l’Homo Habilis.
\subsubsection*{2\quad L’\textit{Homo Erectus}}
L’\textit{Homo Erectus} (« homme dressé, droit ») aurait vécu entre environ 2 millions d’années et 100 000 ans avant notre ère. L’Homme de Pékin par exemple, appartient à l’espèce\textit{Homo Erectus}.

Pour beaucoup de scientifiques, l’Homo Erectus devait posséder le langage car sa production d’outils en silex nécessitait une maîtrise et une technique qui ne pouvait se transmettre et s’enseigner qu’avec le langage. Son organisation sociale (il vivait en campement) va aussi dans ce sens : c’était en quelque sorte un « animal politique ». Mais il y a un autre indice physiologique qui laisse à penser qu’il possédait un langage élaboré : en se tenant debout sur ses deux jambes, en plus de libérer ses mains et d’améliorer ses capacités manufacturières, l’Homo Erectus a surtout appris à mieux contrôler son souffle, tandis que dans le même temps, son thorax s’est élargi pour renforcer la respiration. Or, comme nous le verrons dans le chapitre 7 (Phonétique articulatoire), l’air est indispensable à la phonation, et la source d’air nécessaire à la vibration des cordes vocales vient des poumons..
\subsubsection*{3\quad L’\textit{Homo Sapiens}}
\textit{Sapiens} est un adjectif latin signifiant « intelligent, sage, raisonnable ». Les premiers représentants d’\textit{Homo Sapiens} sont apparus en Afrique entre 200 000 et 100 000 ans avant notre ère.\ednote{这个时间区间过晚。作者原注 9 已提到的 Jebel Irhoud 化石，在 Hublin 等人 2017 年论文中与约 $315\pm34$ 千年前的年代联系，早于此处的 20 万—10 万年。可更正为：至少约 30 万年前，非洲已出现被该研究归入 Homo sapiens 早期阶段的化石；不能据本句时间区间界定首次出现。\ecite{4}} L’Homo Sapiens pourrait être la cause de l’extinction de l’Homme de Néandertal. Parmi les très nombreuses raisons évoquées avec prudence, on cite ses facultés langagières développées, qui lui auraient permis de prendre l’ascendant sur son rival néandertalien, considéré comme fruste, attardé et encore proche du singe.

Et en effet, il apparaît que l’augmentation considérable de la masse cérébrale chez l’Homo Sapiens a été un point-clé dans la maturation de son langage\origfoot{9}{En 2017, on a retrouvé au Maroc sur le site de Jbel Irhoud les plus anciens fossiles d’Homo Sapiens. L’étude des restes de crânes de cinq individus a montré l’expansion de la boîte crânienne et l’augmentation considérable du volume du cerveau de l’Homo Sapiens : de 1100 cm\textsuperscript{3} pour les premiers sapiens à 1 650 cm\textsuperscript{3} il y a 40 000 ans, sachant que notre capacité est en moyenne de 1400 cm\textsuperscript{3}.}. L’Homo Sapiens se livrait aussi à des rites funéraires, autant d’activités impliquant l’existence d’une pensée symbolique et d’un véritable langage.

Il existe actuellement deux principaux scénarios d’apparition de l’Homo Sapiens : un scénario selon lequel il vient d’une source unique, l’Afrique : c’est la monogenèse. Selon l’autre scénario, l’Homo Sapiens est apparu dans plusieurs endroits distincts : c’est la polygenèse.
\subsection[Le point de vue de la génétique]{Le point de vue de la génétique\origfoot{10}{génétique : Sous-discipline de la biologie, la génétique est la science qui étudie l’hérédité et les gènes.}}
En 2011, dans la revue \textit{Science}, \textbf{Quentin Atkinson}, linguiste et anthropologue à l’Université \origcont{33}{019}d’Auckland (Nouvelle-Zélande), a proposé le sud-ouest de l’Afrique comme berceau d’origine de l’ensemble des langues humaines. Fondant son étude sur une analyse comparative du nombre de phonèmes (les sons distincts) présents dans environ 500 langues actuelles, il avait constaté que leur nombre décroît au fur et à mesure que l’on s’éloigne de l’Afrique occidentale (où les phonèmes sont les plus nombreux) pour aller vers les régions du monde les plus tardivement colonisées par l’Homo Sapiens (notamment l’Amérique et l’Océanie, où ces sons élémentaires sont les moins nombreux). Il en a conclu que cette abondance de phonèmes en Afrique de l’ouest faisait de cette partie du continent le creuset de l’élaboration du langage articulé, lequel s’était de plus en plus appauvri en s’éloignant de sa source.
\subsubsection*{1\quad Diversité génétique et diversité linguistique}
Biologiste de formation, Atkinson s’est appuyé sur un parallèle audacieux entre les processus linguistiques et les processus génétiques. Dans le domaine de la génétique des populations, les biologistes observent en effet une diversité génétique maximale en Afrique, considérée de façon consensuelle comme le lieu de naissance de notre espèce, diversité qui diminue en Europe et en Asie pour devenir minimale en Amérique et en Océanie, terres plus récemment occupées : les hommes, bougeant par petits groupes d’individus proches, n’amènent à chaque migration qu’une partie de leur patrimoine génétique initial, qui se réduit de plus en plus à chaque étape. Atkinson avait imaginé qu’il en était de même pour la richesse linguistique.
\subsubsection*{2\quad Les limites de l’analogie}
Oui mais voilà : une année plus tard, des chercheurs européens ont publié la même revue une remise en cause de la méthodologie de l’étude d’Atkinson. Pour éprouver ses résultats, \textbf{Michael Cysouw} et \textbf{Steven Moran}, linguistes à l’Université Louis-et-Maximilien de Munich (Allemagne), et \textbf{Dan Dediu}, de l’Institut Max Planck de psycholinguistique de Nimègue (Pays-Bas), l’ont appliquée sur d’autres critères linguistiques (ex. l’utilisation de la forme passive). Les résultats ne pointent pas dans la même direction, et selon les cas, désignent plutôt comme berceau du langage l’Afrique orientale, le Caucase ou d’autres régions totalement différentes.

Malgré cette réserve d’ordre méthodologique, d’autres recherches en paléolinguistique ont identifié un fonds de mots communs toutes les langues terrestres écrites au début du 21\textsuperscript{e} siècle, ce qui tend à favoriser le scénario de la monogenèse. On peut donc penser qu’au moment de la « sortie d’Afrique » et avant de se répandre sur la planète, l’Homo Sapiens possédait déjà le langage. L’Afrique serait bien le berceau du langage.
\subsection{Le point de vue de la paléogénétique}
La paléogénétique est une science récente qui s’intéresse à la récupération et à l’analyse de l’ADN\origfoot{11}{L’ADN (L’acide désoxyribonucléique) est une macromolécule biologique présente dans toutes les cellules. Il contient toute l’information génétique (appelée génome) permettant le développement, le fonctionnement et la reproduction des êtres vivants.} des organismes fossiles. Certains paléogénéticiens considèrent que c’est lors du passage à l’Homo \origcont{34}{020}Sapiens que sont apparues les aires de Broca et de Wernicke\origfoot{12}{Les aires de Broca et de Wernicke, ainsi que la découverte de FOXP2 seront présentées en détail dans notre chapitre 15 : Biolinguistique et psycholinguistique.} (zones du cerveau responsable du langage) suite à la mutation d’un gène, le gène FOXP2, appelé « gène de la parole ». C’est cette mutation qui aurait donné la capacité à l’homme de passer des simples mots à la syntaxe, et donc au langage articulé.\ednote{这里把一个相关基因写成了语法和脑区起源的单一充分原因，证据不足。Lai 等人（2001）发现的是 FOXP2 变异与严重言语和语言障碍的关联，并未证明某一次变异产生 Broca、Wernicke 区，或独自使人类获得句法。宜改为：FOXP2 参与言语和语言相关的神经发育过程；不能称为单独决定语言能力的“语言基因”。\ecite{5}}
\subsubsection*{1\quad La thèse de Philip Lieberman}
L’aptitude physique au langage articulé des Néandertaliens a longtemps été controversée. La paléogénétique ayant d’abord calculé que notre gène FOXP2 serait vieux de 100 000 à 200 000 ans, on en avait déduit que l’Homme de Néanderthal, déjà présenté comme très inférieur à l’Homo Sapiens, ne disposait pas du langage. Par ailleurs, certains scientifiques américains, comme \textbf{Philip Lieberman} (né en 1934), auteur de \textit{On the Origins of Language} (1975) et \textbf{Edmund Crelin} (1923-2004), professeur d’anatomie à Yale, ont nié la possibilité que le langage articulé ait pu exister chez les prédécesseurs de l’Homo Sapiens. Selon eux, un larynx en position basse est nécessaire pour produire un langage articulé. Par ailleurs, le Néandertalien ne disposait pas d’un pharynx de taille suffisante pour produire tous les sons que l’on observe dans les langues du monde.
\subsubsection*{2\quad L’homme de Néandertal : un grand bavard ?}
La thèse sur l’âge du gène FOXP2 a été révisée en octobre 2007 lorsqu’une équipe de scientifiques dirigée par le paléogénéticien \textbf{Johannes Krause} (né en 1980), du Max Planck Institute for Evolutionary Anthropology de Leipzig, a étudié le gène FOXP2 provenant des restes de deux Néandertaliens, et n’a trouvé aucune différence avec notre gène FOXP2. Le néandertalien disposait donc aussi des mutations de FOXP2 rendant possible la parole.\ednote{Krause 等人（2007）的主要结果是尼安德特人与现代人共享所考察的两个衍生氨基酸变化，并非证明整个 FOXP2 基因所有位点都完全相同，更非证明现代人式语言能力已经存在。此处应把“没有任何差别”限定为该研究考察的变异，把后一句的确定性结论改为对相关生物学条件的证据。\ecite{6}} Quant à théorie de Lieberman, malgré de nombreuses critiques, elle s’est largement diffusée pendant une trentaine d’années, avant d’être finalement infirmée par toute une série de découvertes récentes :
\begin{itemize}
\item Certains fossiles de Néandertaliens montrent des larynx en position basse, donc permettant la phonation.
\item On a découvert en 1983 en Israël un os hyoïde daté de 60 000 ans avant J.-C. sur un squelette de Néandertal. En 2008, des os hyoïdes appartenant à des pré-néandertaliens (datés d’au moins 530 000 ans avant J.-C.) ont été découverts sur le site préhistorique de la Sima de los Huesosen Espagne.\ednote{底本地名与介词黏连：« Huesosen Espagne » 应为 « Huesos en Espagne »。} L’os hyoïde est un petit os qui maintient la base de la langue et dont la morphologie est déterminante pour l’aptitude à l’élocution chez l’être humain : c’est cet os qui permet le mouvement du larynx nécessaire à l’articulation vocale. Ces deux os sont très peu différents de ceux des humains actuels.
\item Enfin, des modélisations informatiques modernes montrent que, même avec un larynx en position haute, les articulateurs (langue, voile du palais, dents et lèvres) mis en jeu lors de la phonation ont un degré de liberté qui permet de produire une large gamme de sons articulés.
\end{itemize}
\origpage{35}{021}
Ainsi, alors que dans les années 80 la majorité des scientifiques, influencés par Lieberman penchaient vers une origine du langage plutôt récente, de l’ordre de 100 000 à 40 000 ans avant notre ère, il est maintenant permis de penser que les origines du langage sont beaucoup plus anciennes. Aujourd’hui, la plupart des spécialistes envisagent l’apparition du langage en deux étapes : une première phase de langage primitif (protolangage) parlé par l’\textit{Homo Erectus}, suivie par le langage complexe avec \textit{Homo Sapiens}.

Des études précises sur l’appareil vocal ont montré que l’\textit{Homo Erectus} disposait d’une palette de sons identique à celle d’un enfant de deux ans. Or, avec quelques phonèmes, on peut déjà produire un vocabulaire assez varié. Par ailleurs, des recherches sur l’anatomie du cerveau suggèrent également de reculer les dates d’apparition du langage. En effet, depuis 2,5 millions d’années jusqu’à aujourd’hui, on constate chez tous les Homo une augmentation continue de la taille des zones du cerveau où sont centralisées les activités de conceptualisation, de planification et de langage.

Les dernières études sur les aptitudes anatomiques des premiers hominidés repoussent donc les origines du langage à 2 millions d’années.\ednote{本段把发声的解剖条件直接等同于语言已经出现，推论过强。即使证明某古人类能产生若干语音，也不能仅由此推出其具备语法、语义和交际系统，更不能由此把“语言诞生于 200 万年前”写成已经确立的年代。这里宜表述为有关早期发声条件或原始语言的假说；本文所给材料不足以判定具体起源日期。}

C’est ainsi que se termine notre premier chapitre. Dans le chapitre suivant, « Langage humain et communication animale », nous essaierons de répondre à une autre question longtemps débattue par philosophes et scientifiques : le langage est-il le propre de l’homme ?
\Needspace{11\baselineskip}\section*{Exercices}\phantomsection\addcontentsline{toc}{section}{Exercices}\markright{Exercices}
\noindent\textbf{I.\quad Dans cet extrait de \textit{Philosophie de l’Esprit} (1830), de quelle manière Georg Friedrich Hegel (1770-1831) relie-t-il les mots à la pensée ? Quel lien pouvez-vous établir entre cette conception et celle de Rousseau ? Par ailleurs, comment comprenez-vous la notion d’ineffable ? Pensez-vous que les mots ne peuvent pas tout dire ?}
\begin{adjustwidth}{4mm}{4mm}\small
C’est dans les mots que nous pensons. Nous n’avons conscience de nos pensées déterminées et réelles que lorsque nous leur donnons la forme objective, que nous les différencions de notre intériorité, et par suite, nous les marquons d’une forme externe, mais d’une forme qui contient aussi le caractère de l’activité interne la plus haute. C’est le son articulé, le mot, qui seul nous offre une existence où l’externe et l’interne sont si intimement unis. Par conséquent, vouloir penser sans les mots, c’est une tentative insensée. [...] Et il est également absurde de considérer comme un désavantage et comme un défaut de la pensée cette nécessité qui lie celle-ci au mot. On croit ordinairement, il est vrai, que ce qu’il y a de plus haut, c’est l’ineffable. Mais c’est là une opinion superficielle et sans fondement ; car, en réalité, l’ineffable, c’est la pensée obscure, la pensée à l’état de fermentation, et qui ne devient claire que lorsqu’elle trouve le mot. Ainsi le mot donne à la pensée son existence la plus haute et la plus vraie.
\end{adjustwidth}
\origpage{36}{022}
\noindent\textbf{II.\quad Dans cet extrait du \textit{Traité sur l’origine des langues} (1771), Johann Gottfried von Herder (1744-1803), oppose les forces et les faiblesses des hommes et des animaux : retrouvez-les. Par ailleurs, comment comprenez-vous l’idée de perfectionnement ? Pensez-vous comme Herder que seul l’homme peut se perfectionner ?}
\begin{adjustwidth}{4mm}{4mm}\small
Si l’homme possède des forces plastiques qui ne se limitent pas à construire une cellule de miel ou à tisser une toile d’araignée, des forces qui, en ce domaine, sont inférieures aux facultés industrieuses des bêtes, celles-ci ont aussi une tout autre perspective. L’homme n’est pas condamné à œuvrer toujours sur une seule chose ; il a le champ libre pour s’exercer sur beaucoup, et se perfectionner sans cesse. Ses idées ne sont pas un pur produit de la nature, mais peuvent par conséquent devenir également son œuvre propre.
\end{adjustwidth}
\section*{Pour aller plus loin}\phantomsection\addcontentsline{toc}{section}{Pour aller plus loin}\markright{Pour aller plus loin}
\begingroup\small\setlength{\parskip}{.15em}\setlength{\parindent}{-4mm}\leftskip=4mm
COPPENS (Yves) et PICQ (Pascal). \textit{Aux origines de l’humanité}. Tome 1, De l’apparition de la vie à l’homme modern, Tome 2, Le propre de l’homme, Fayard, 2001.

DESSALLES (Jean-Louis). \textit{Aux origines du langage}. Une histoire naturelle de la parole, Hermès, 2000.

HERDER (Johann-Gottfried). \textit{Traité sur l’origine des langues}. Paris, Allia, 2010.

HOMBERT (Jean-Marie) et Al. \textit{Aux origines des langues et du langage}. Paris, Fayard, 2005.

LIBERMAN (Philip). \textit{On the Origins of Language}. New York, Macmillan, 1975.\ednote{原书目作 « LIBERMAN »，与本章正文反复出现的 « Lieberman » 不一致；应更正为 « LIEBERMAN (Philip) »。}

MACAGNO (Gilles). \textit{La longue marche d’Homo Sapiens}. La fabuleuse histoire du bipède, Ellipses, 2005.

PIATELLI-PALMARINI (Massimo). \textit{Théories du langage, théories de l’apprentissage}. Seuil, 1979.

RULEN (Merritt). \textit{The Origin of language}. Tracing the evolution of the mother tongue, New York, Wiley, 1996.
\par\endgroup
\clearpage
\origpage{37}{023}
\section*{Glossaire}\phantomsection\addcontentsline{toc}{section}{Glossaire}\markright{Glossaire}
\gloss{analogie\quad 类推法}{语言学习的一个过程，学习者根据他所知的其他形式的模式来建构位置形式。\ednote{底本作“位置形式”。依“根据已知形式推知新形式”的上下文，此处疑为“未知形式”之误，属于拟校，尚未获得编者所据词典的对应条目确证；不将拟校静默写回正文。}}
\gloss{discours\quad 话语；语篇}{语言运用实例的通称，即作为交际行为的结果的语言。在后现代主义和批评话语分析中，话语不但被用来指任何交谈，而且指蕴含在交谈中的意义和价值。}
\gloss{langage\quad 语言\textsuperscript{1}}{指人类的交际系统，由有组织的语音系列（或其书面形式）构成。这些语音系列组成更大的单位，如语素、单词、句子、话语等。}
\gloss{langue\quad 语言\textsuperscript{2}}{任何特定的人类交际体系，例如，汉语、法语、印地语。}
\gloss{langue maternelle\quad 母语}{通常指在家里学会的第一语言。}
\gloss{langue sémitique\quad 闪语}{属于阿非罗－亚细亚语系，该语系是分布于北非和西南亚的一个主要语系。闪语族主要包括阿拉伯语、希伯来语、阿拉米语、马耳他语和阿姆哈拉语等。}
\gloss{linguistique\quad 语言学}{把语言作为人类交际系统来研究的一门学科，研究对象包括语音系统、句子结构、语言和认知关系、意义系统以及语言和社会因素等。}
\gloss{parole\quad 言语}{交际时作为功能性单位的语句，语句包含字面意义和言外意义。}
\gloss{philosophie\quad 哲学}{关于世界观、价值观、方法论的学说。是在具体各门科学知识的基础上形成的，具有概括性、抽象性、反思性、普遍性的特点。}
\gloss{phonation\quad 发音；发声}{发出语音或乐音，也泛指发出声音。}
